%HOMEWORK template for M 325K
\documentclass{article}
\headheight=7pt         \topmargin=14pt
\textheight=574pt       \textwidth=445pt
\oddsidemargin=18pt     \evensidemargin=18pt

%Begin package import

\usepackage{amsmath, amssymb, amsthm}
\usepackage{mathtools}
\usepackage{pdfpages}
\usepackage{tikz-cd}

%End package import
%Begin custom commands

\newcommand*{\T}{\mathcal{T}}
\newcommand*{\B}{\mathcal{B}}
\newcommand*{\U}{\mathcal{U}}
\newcommand*{\cl}{\mathrm{cl}}
\newcommand*{\subeq}{\subseteq}
\newcommand*{\empset}{\emptyset}
\newcommand{\C}{\mathbb{C}}
\newcommand{\R}{\mathbb{R}}
\newcommand{\Q}{\mathbb{Q}}
\newcommand{\PR}{\mathbb{P}}
\newcommand{\Z}{\mathbb{Z}}
\newcommand{\N}{\mathbb{N}}
\newcommand{\F}{\mathcal{F}}
\newcommand{\abs}[1]{\left\lvert #1\right\rvert}
\newcommand{\RM}[1]{\mathrm{#1}}
\newcommand{\BF}[1]{\textbf{#1}}
\newcommand{\e}{\epsilon}
\newcommand{\ceil}[1]{\lceil{#1}\rceil}
\newcommand{\floor}[1]{\lfloor{#1}\rfloor}
\newcommand{\rarr}[0]{\rightarrow}
\newcommand{\IM}[1]{\text{Im}(#1)}
\newcommand{\Hom}[0]{\text{Hom}}
\newcommand{\Pow}[1]{\text{Pow}(#1)}
\newcommand{\angles}[1]{\langle #1 \rangle}


%End custom commands
%Begin custom theorems

\newtheorem{innercustomgeneric}{\customgenericname}
\providecommand{\customgenericname}{}
\newcommand{\newcustomtheorem}[2]{%
\newenvironment{#1}[1]
{%
\renewcommand\customgenericname{#2}%
\renewcommand\theinnercustomgeneric{##1}%
\innercustomgeneric%
}
{\endinnercustomgeneric}
}

\newcustomtheorem{THRM}{Theorem}
\newcustomtheorem{LEM}{Lemma}
\newcustomtheorem{EX}{Exercise}
\newcustomtheorem{COR}{Corollary}
\newcustomtheorem{PROB}{Problem}
\newcustomtheorem{claim}{Claim}

\newenvironment{solution}
{\renewcommand\qedsymbol{$\blacksquare$}\begin{proof}[Solution]}
{\end{proof}}

\newenvironment{definition}
{\renewcommand\qedsymbol{$\blacksquare$}\begin{proof}[Definition]}
{\end{proof}}

%End custom theorems
\begin{document}

\title{Exercise 04}
\author{Arham Lodha}%put your name here
\date{March 12, 2025}%enter the due date here
\maketitle

\begin{PROB}{1a}\label{Problem 1a.}
\end{PROB}
\begin{solution}
    \begin{enumerate}
        \item False
        \item False
        \item True
        \item True
        \item True
    \end{enumerate}
\end{solution}

\begin{PROB}{1b}\label{Problem 1b.}
\end{PROB}
\begin{solution}
    \begin{enumerate}
        \item True
        \item False
        \item True
        \item False
        \item True
    \end{enumerate}
    
\end{solution}

\bigskip

\begin{PROB}{2a}\label{Problem 2a.}
\end{PROB}
\begin{proof}
    $\forall x \in \Z$ and let $\alpha = \frac{1}{2}$. We want to show that $E_x(V_x) = \infty$. Note $p_{x, x}^{2n} = 2^{-2n}{2n \choose n} \sim \frac{c}{\sqrt{n}}$ as $n \to \infty$ as given.   
    
    \begin{align*}
        E_x(V_x) &= \sum_{n = 1}^{\infty}p_{0, 0}^{(n)} \\
        &= \sum_{n = 1}^{\infty}p_{0, 0}^{(2n)} \\
        &= \sum_{n = 1}^{\infty}{2n \choose n}\left(\frac{1}{2}\right)^{2n} \\
        &\geq \sum_{n = 1}^{\infty} \frac{c}{2\sqrt{n}} = \infty
    \end{align*}

    Thus $E_x(V_x) = \infty$. Thus $x$ is recurrent.  
\end{proof}
\begin{PROB}{2b}\label{Problem 2b.}
\end{PROB}
\begin{proof}
    Let $\left\{ Y_i \right\}_{i \geq 1}$ be a sequence of random variables where $\PR(Y_i = 1) = \alpha$ and $\PR(Y_i = -1) = 1- \alpha$. Thus $X_n = \sum_{m = 1}^{n}Y_m$.
    
    By the Strong Law of Large Numbers, $\frac{X_n}{n} \to E[Y_m] = \alpha - 1  + \alpha = 2 \alpha - 1$ almost surely. For $\alpha < \frac{1}{2}$: $\frac{X_n}{n}$ approaches a value less than 0. Thus $X_n \to -\infty$. Similarly for $\alpha > \frac{1}{2}$: $\frac{X_n}{n}$ approaches a value greater than 0. Thus $X_n \to \infty$. Thus in either case: $x$ is a transient state.       
    
\end{proof}

\begin{PROB}{2c}\label{Problem 2c.}
    
\end{PROB}
\begin{proof}
    Let $Y_n = \abs{X_n}$. 
    
    \begin{enumerate}
        \item For $\alpha = \frac{1}{2}$, $X_n$ is recurrent at 0. Thus $Y_n$ is recurrent at 0. Let $Z_n$ be such a reflected walk.    
        \item For $\alpha < \frac{1}{2}$: Couple $Y_n$ and $Z_n$ in the following way: 
    
        \begin{enumerate} 
            \item If $Y_n = 0$: Move right. Otherwise 
            \item If $Z_n$ moves left then $Y_n$ moves left. 
            \item If $Z_n$ moves right then $Y_n$  we pick from the leftover portion of 
            Y's distribution for the “right step.” Since
            $\alpha < \frac{1}{2}$, this leftover portion is enough to ensure we can couple them consistently.     
        \end{enumerate}

        Let $T_k$ be the kth visiting time of $Z$ on 0. We want to show for all $k \geq 1$, $\exists n \in [T_k, T_{k + 1}] \ni Y_n = 0$. Let $l_z, l_y, r_z, r_y$, be the left steps of $Z$ and $Y$  and the right steps of $Z$ and $Y$.
        
        Assume the contrary $\forall n \in [T_k, T_{k + 1}]$, $Y_n \neq 0$. 
        In the interval, $[T_k, T_{k + 1}]$, $l_z = r_z$. Also note because of the coupling, $l_y \geq l_z$. Since $l_z + r_z = T_{k + 1} - T_{k} = l_y + r_y$, at $Y_{T_{k + 1}} \leq Z_{T_{k + 1}} = 0$. However $Y_{T_{k + 1}} \geq 0 \implies Y_{T_{k + 1}} = 0$, which is a contradiction. 
        
        Thus $\exists n \in [T_k, T_{k + 1}] \ni Y_n = 0$. Thus $Y$ visits $0$ infinitely often since $Z$ visits $0$ infinitely often. Thus $0$ is recurrent for $Y$.       
    \end{enumerate}
    
\end{proof}

\bigskip


\begin{PROB}{3}\label{Problem 3.}
    
\end{PROB}
\begin{proof}
\end{proof}
    Let $Y_n$ be a SRW on a d-regular tree. 
    $\forall x \in V$. Let $D_n := \abs{Y_n - x}$ or distance from $Y_n$ to x. $D_n$ is a biased reflected random walk where $\alpha = \frac{d - 1}{d}$. Thus $D_n$ isn't transient which means $Y_n$  isn't transient on $x$. 
\bigskip

\bigskip

\begin{PROB}{4a}\label{Problem 4a.}
\end{PROB}
\begin{proof}
    $\forall (x, y) \in \Z^2 \setminus \left\{ (0, 0) \right\}$, $\PR_{(x, y)}(H_{(0, 0)} = \abs{x} + \abs{y}) = 1$. 
    
    Thus, 
    
    \begin{align*}
        \PR_{(0, 0)}(H_{(0,0)} < \infty) &= \sum_{(x, y) \in \Z^2} \PR_{(0, 0)}(X_1 = (x, y)) \PR_{(x, y)}(H_{(0, 0)} < \infty \mid X_1 = (x, y)) \\
        &= \mu(0,0) + \sum_{(x, y) \in \Z^2 \setminus {(0, 0)}} \mu(x, y) \PR_{(x, y)}(H_{(0, 0)} < \infty) \\
        &= \mu(0,0) + \sum_{(x, y) \in \Z^2 \setminus {(0, 0)}} \mu(x, y) 
        &= 1
    \end{align*}
    
    Thus $(0 ,0)$ is recurrent.  
\end{proof}

\begin{PROB}{4b}\label{Problem 4b.}
\end{PROB}
\begin{proof}
    \begin{align*}
        E_0(H_0) &= \PR_{(0, 0)}(X_1 = (0, 0)) + \sum_{(x, y) \in \Z^2 \setminus {0, 0}} \PR_{(0, 0)}(X_1 = (x, y))\left(E_{(x, y)}(H_0) + 1\right) \\
        &= 
        \mu(0, 0) + \sum_{(x, y) \in \Z^2 / \left\{ (0, 0) \right\}} \mu(x, y)(\abs{x} + \abs{y} + 1) \\
        &= 1 + \sum_{(x, y) \in \Z^2 / \left\{ (0, 0) \right\}} \mu(x, y)(\abs{x} + \abs{y}) \\
        &= 1 + E_\mu[\abs{x} + \abs{y}]
    \end{align*}

    $(0,0)$ is positive recurrent if $E_\mu[\abs{x} + \abs{y}]$ is finite and null recurrent otherwise.   
\end{proof}

\begin{PROB}{4c}\label{Problem 4c.}
\end{PROB}
\begin{proof}
    \begin{enumerate}
        \item Let $Y_n = \abs{\abs{X_n}}_\infty$.
        
        \begin{align*}
            \PR_0(Y_2 = 2 \mid Y_1 = \abs{\abs{(2, 2)}}_{\infty} = 2) &= 1
        \end{align*}

        and \begin{align*}
            \PR(Y_2 = 2 \mid Y_1 = 2) &= \frac{1}{2}
        \end{align*}

        Thus $Y_n$ is not a markov processs since $\PR(Y_2 = 2 \mid Y_1 = 2) \neq \PR_0(Y_2 = 2 \mid Y_1 = \abs{\abs{(2, 2)}}_{\infty} = 2)$  

        \item Let $Y_n = \abs{\abs{X_n}}_{1}$. From any state with $L^1$ norm $k$ for $k \neq 0$ , the next norm is $k - 1$ deterministically. For $k = 0$, 
        
        $p_{k, n} = \sum_{\abs{x } + \abs{y} = n} \mu(x, y)$. 
        

        The transitions depend only on the current norm, satisfying the Markov property

        \item The projection onto the x-coordinate is not Markov. The transition of the x-coordinate depends on the underlying y-coordinate, which is not captured by the projection. Different underlying states with the same x-coordinate can lead to different future x-coordinates, violating the Markov property
    \end{enumerate}
    
\end{proof}


\end{document}